\chapter{What the linear commutator sews}\label{ch:11}
Two detectors, one on either side of a cut, can be used in two different
experiments. In the first, a source at one detector is varied and the
response at the other is recorded. In the second, the readings fluctuate in
repeated preparations of the same state and their correlation is measured.
For a linear bosonic field the first experiment fixes a commutator. It does
not in general predict the second experiment. Keeping the two questions
separate lets us quantize a sewn system even when neither a vacuum nor a
factorization of its Hilbert space has been specified.

\section{From physical linear observables to canonical commutators}
Fix the geometric identification, exterior boundary condition and classical
sector, and linearize about a solution in that sector. Tangents satisfying
the linearized equations, constraints and boundary conditions carry the
presymplectic form $\omega=\delta\theta$ of Chapter~\ref{ch:05}, including
its corner term. A proper gauge tangent pairs to zero with every allowed
tangent. A transformation with a variable surface charge generally does
not. Write $\widetilde{\mathcal S}$ for these constrained solution
tangents and $\mathcal G_{\rm proper}$ for their proper gauge directions.
The physical tangent space is $\mathcal S_{\rm phys}=
\widetilde{\mathcal S}/\mathcal G_{\rm proper}$. To call it symplectic
requires that on the specified sector no further radical of $\omega$
remains. Dividing by every transformation bearing a gauge parameter would
not establish that. A central charge fixed in the sector is a number; a
physical zero mode with a conjugate variable is still a canonical pair.
This distinction matters especially at a cut, where a charged transformation
on one piece can become a redundancy only after joining the pieces.

Here is a test-function realization that fixes the signs. For the real
scalar take $P=\Box-m^2-\xi R$, the Hessian with the sign of the action in
Chapter~\ref{ch:04}; $R$ here denotes scalar curvature. Use the metric
volume $dV$. Write $G_R,G_A$ for the retarded and advanced inverses
($G^R,G^A$ in Chapter~\ref{ch:06}) of this $P$, with the same physical
boundary condition and adjoint to one another. Put
\[
 E=G_R-G_A,\qquad F_f(\phi)=\int_M f\phi\,dV,\qquad
 E(f,h)=\int_M f(Eh)\,dV .
\]
Real sources $f,h$ are smooth and compactly supported in the interior. On
an ordinary globally hyperbolic spacetime without boundary, the label space
is $V=C_c^\infty(M;\mathbb R)/P C_c^\infty(M;\mathbb R)$: $F_{Pf}$
vanishes on solutions, and $E$ identifies this quotient with the real
solutions of spacelike-compact support \cite{WaldQFT}. At a physical wall
the test domain and Green identity must implement the wall condition. For
gauge fields one uses admissible linear observables that annihilate proper
gauge variations, modulo the equations and null observables. A propagator
of a gauge-fixed equation on all components is not by itself the physical
label space.

With $\theta=p\,\delta q$ on a canonical pair, the convention is
$\omega=\delta p\wedge\delta q$. Green's identity for the scalar gives
$\delta F_f(Y)=\omega(Y,Ef)$: $Ef$ is its Hamiltonian tangent. In
particular, the antisymmetric form on \emph{observable labels} is
\begin{equation}
 \sigma([f],[h])=\{F_f,F_h\}_{\mathrm P}
 =\delta F_f(Eh)=E(f,h)=-\omega(Ef,Eh).
 \label{eq:11-label-sign}
\end{equation}
The last minus sign trades the order of the arguments in the form on
solutions; it is not a change from retarded minus advanced. One check is a
spatial mode for which $P=-(\partial_t^2+\nu^2)$. Its causal kernel is
$-\sin[\nu(t-t')]/\nu$. Differentiating with respect to $t'$ at $t'=t$
gives $+1$, the bracket of displacement with its momentum.

\begin{definition}
For a real symplectic label space $(V,\sigma)$ the Weyl form of the
bosonic canonical commutation relations (CCR) defines the universal
algebra generated by unitaries $\mathcal W(v)$, $v\in V$, with
$\mathcal W(0)=\mathbf 1$ and
\begin{equation}
 \mathcal W(v)^*=\mathcal W(-v),\qquad
 \mathcal W(v)\mathcal W(w)=
 e^{-i\hbar\sigma(v,w)/2}\,\mathcal W(v+w).
 \label{eq:11-weyl}
\end{equation}
A regular representation is one in which $s\mapsto\mathcal W(sv)$ is
strongly continuous for every $v$. Its generator is written
$\mathcal W(sv)=e^{is\widehat X_v}$.
\end{definition}

The same relations make sense for a presymplectic space: labels in the
radical give central unitaries until their values are fixed or their labels
removed. The exponential relation avoids multiplying arbitrary unbounded
operators. In the field algebra, or by differentiation on suitable common domains,
$[\widehat X_f,\widehat X_h]=i\hbar E(f,h)\mathbf 1$; subscripts $f,h$
stand for their classes. No representation is needed to specify
\eqref{eq:11-weyl}. We set $\hbar=1$ from now on unless it is displayed.
If a positive-frequency basis is available and chosen, normalize physical
modes with $i\omega(u_I,\overline{u_J})=\delta_{IJ}$ and
$\omega(u_I,u_J)=0$. Then
\[
 \widehat\phi=\sum_I(a_Iu_I+a_I^\dagger\overline{u_I}),\qquad
 [a_I,a_J^\dagger]=\delta_{IJ}\mathbf 1 .
\]
For the oscillator this gives $u=e^{-i\nu t}/\sqrt{2\nu}$ and
$u(t)\bar u(t')-u(t')\bar u(t)=iE(t,t')$. Delta functions and spectral
integrals replace Kronecker symbols and sums for continuous labels. This is
one realization, not an instruction to choose positive frequency separately
on the two sides.

\section{Two routes to an algebra}
Suppose the classical assembly in the chosen sector gives the physical sewn
labels $V_{\rm sew}$ and a comparison $S:V_{\rm sew}\to V_M$ with the
independently specified whole. What licenses quantization of the comparison
is the bracket test
\[
 \sigma_M(Sv,Sw)=\sigma_{\rm sew}(v,w).
\]
For bulk sources this tests the causal propagator, including its blocks
between opposite sides of the cut. For instance, in the massless line
example of Chapter~\ref{ch:08}, a source at $x=-a$ and a detector at
$x=b$, $a,b>0$, are separated by a travel time $a+b$. For $t>t'$ the
transmitted kernel of the Hessian $P=-(\partial_t^2-\partial_x^2)$ gives
\[
 [\widehat\phi(t,b),\widehat\phi(t',-a)]
 =-\tfrac i2\,\Theta(t-t'-a-b)\mathbf1
\]
as a distribution on that time ordering. The minus converts the positive
wave operator used for the pulse to the action Hessian. Two independent
reflectors would give zero even after the first arrival. Matching just the
spaces of homogeneous solutions does not compute this block.

If $S$ is a symplectic isomorphism, $\mathcal W_{\rm sew}(v)\mapsto\mathcal W_M(Sv)$
is a star isomorphism: it preserves the phase in \eqref{eq:11-weyl}, and
$S^{-1}$ supplies the inverse. This is the \emph{sew, then quantize} route.
If the comparison covers only a symplectic subspace, the same argument only
identifies its generated subalgebra.

The converse construction starts with quantum generators of opened pieces.
What it means to impose a matching equation on them is not automatic. Even
for canonical pairs, setting $q_1-q_2=0$ in the full polynomial operator
algebra would set $[q_1-q_2,p_1]=i\mathbf 1$ to zero. The quotient by the
two-sided ideal so obtained would be the zero algebra. One must first find
observables that survive the matching. The following linear case gives a
precise meaning to that instruction.

\begin{proposition}
Let $U$ be a real symplectic space of extended linear labels and let
$I\subset U$ be isotropic: $\sigma(i,j)=0$ for $i,j\in I$. Suppose the
matched presentation is $I^\perp/I$, where $I^\perp$ consists of labels
pairing to zero with $I$. In the Weyl algebra of $I^\perp$, impose
$\mathcal W(i)=\mathbf1$ for $i\in I$. The resulting algebra has precisely
the Weyl relations of $I^\perp/I$ with its induced form. It is a
\emph{symplectic} presentation of the physical labels if the comparison
identifies their form and
$\operatorname{rad}(\sigma|_{I^\perp})=(I^\perp)^\perp=I$. Otherwise
the quotient can still be presymplectic.
\end{proposition}
\begin{proof}
In the indicated subalgebra the constraint unitaries are central. Since
$\mathcal W(v+i)=\mathcal W(v)\mathcal W(i)$ for $v\in I^\perp$, the
class of $\mathcal W(v)$ depends only on $v+I$. The phase of a product
also depends only on the two classes. This gives a map from the Weyl
algebra on the quotient labels; sending each old generator to the Weyl
generator of its class defines the inverse map. The radical of the induced
form is $(I^\perp)^\perp/I$: isotropy gives $I\subset I^\perp$ and hence
$(I^\perp)^\perp\subset I^\perp$. The double-orthogonal condition is
automatic in finite symplectic dimension, not for every infinite-dimensional
choice of subspace. Residual central Weyl generators must not silently be
set to one.
\end{proof}

For $I$ generated by $q_1-q_2$, the surviving total momentum is
$p_1+p_2$, not either $p_1$ or $p_2$ alone. With
$Q=(q_1+q_2)/2$ one finds $[Q,p_1+p_2]=i\mathbf1$: no factor of two is
lost when the displacements are identified. The proposition treats a
linear first-class constraint in its zero sector. A nonzero fixed charge
requires its specified central character; noncommuting (second-class)
constraints generally change the bracket of the survivors. A retarded
boundary-history problem need not admit this constraint presentation at
all. In particular, the direct sum of the label spaces of two independently
\emph{reflecting} theories carries their reflecting commutators. It does
not acquire a transmitting commutator by a choice of quantum state. A
quantize-first route must actually construct the physical label space and
prove the bracket comparison on its claimed domain.

\section{Extending correlations, not just commutators}
A state $\varpi$ is a normalized positive linear functional on the algebra;
positivity means $\varpi(A^*A)\ge0$. A centered quasifree (Gaussian)
state is specified by a real symmetric covariance $\mu$:
\begin{align}
 \varpi(\mathcal W(v))&=e^{-\mu(v,v)/2},\nonumber\\
 W(v,w):=\varpi(\widehat X_v\widehat X_w)
 &=\mu(v,w)+\tfrac i2\sigma(v,w) .\label{eq:11-covariance}
\end{align}
A mean would add its linear phase to the first formula and the product of
the means to the second. Higher centered moments are sums over pairings of
$W$, in the order of the fields. The necessary and sufficient Gaussian
positivity condition is that the Hermitian form $\mu+i\sigma/2$ be positive
semidefinite on the complexification of $V$ \cite{WaldQFT,Haag}. Indeed,
its quadratic form on $z=\sum c_j v_j$ is the expectation of
$(\sum c_j\widehat X_{v_j})^*(\sum c_j\widehat X_{v_j})$. In particular
\[
 \mu(v,v)\mu(w,w)\ \ge\ \mu(v,w)^2+\tfrac14\sigma(v,w)^2 .
\]
With $\hbar$ restored, replace $\sigma$ by $\hbar\sigma$ in these state
formulas. The commutator fixes the imaginary part; even a completely known
bracket therefore leaves a positivity problem for the real part.

For example, suppose two admissible regional subspaces have been embedded
in $V_M$. Seek a single covariance on $V_M$ whose restrictions are the
given $\mu_1,\mu_2$. They must agree on common data. In a presentation by
regional generators, the trial cross block $C(v_1,v_2)$ must make the
whole symmetric form descend through every relation: a label representing
zero must pair to zero with \emph{all} labels. For finite families of test
labels on the two sides the positivity test is the matrix inequality
\[
 \begin{pmatrix}
   \mu_1+i\sigma_{11}/2 & C+i\sigma_{12}/2\\
   (C+i\sigma_{12}/2)^* & \mu_2+i\sigma_{22}/2
 \end{pmatrix}\ \ge 0 .
\]
Here the $\sigma_{ab}$ are the blocks of the already sewn commutator, and
$*$ is conjugate transpose; the test must hold for every such family. If
the labels span only part of $V_M$, a positive extension to the rest is
still needed. For a field state the resulting two-point functional must
also be a distribution with the equations, boundary conditions and
short-distance behavior required of the whole theory. Equality of two
marginals alone performs none of these tests.

Take two commuting oscillator pairs ordered as $(Q_1,P_1,Q_2,P_2)$.
Write $J_2=\left(\begin{smallmatrix}0&1\\-1&0\end{smallmatrix}\right)$
and $1_2$ for the identity matrix. Consider
\[
 \sigma=J_2\oplus J_2,\qquad
 \mu_c=\begin{pmatrix}\nu 1_2&c1_2\\c1_2&\nu 1_2\end{pmatrix}.
\]
The sum and difference coordinates $(Q_1\pm Q_2)/\sqrt2$ and their
similarly normalized momenta give two canonical pairs with covariances
$(\nu+c)1_2$ and $(\nu-c)1_2$. Positivity holds precisely when
$\nu-|c|\ge1/2$. Thus $\nu=1$ admits both $c=0$ and $c=1/4$, with
identical individual Gaussian states and identical commutators, but
$\varpi(Q_1Q_2)=c$. A putative extension with $c=3/4$ fails the uncertainty
test in the difference pair. This finite calculation proves neither
nonuniqueness nor existence for every spacetime cover. It isolates exactly
which matrix entry the two regional commutator experiments did not measure.

\section{A sharp spatial cut and the vacuum}
There is a local version of this diagnostic that directly tests a proposed
product vacuum. In four-dimensional Minkowski space the massless scalar
vacuum has, at equal times and distinct spatial points,
\[
 W((t,\mathbf x),(t,\mathbf y))=
 \frac{1}{4\pi^2|\mathbf x-\mathbf y|^2},\qquad
 E((t,\mathbf x),(t,\mathbf y))=0 .
\]
The formulas may be smeared over two disjoint small balls on opposite
sides of a spatial plane. The first pairing is positive for nonzero
nonnegative test functions; the second is zero. The product of the two
centered marginal states would predict zero for both. As the balls approach
the same point from opposite sides, the missing correlation is the leading
short-distance singularity of the field, not an arbitrary convention at
one point. The Hadamard condition requires this local singularity on the
\emph{uncut} spacetime. It does not pick the vacuum: adding a smooth,
positive kernel $c\,u(x)u(y)$ with $c\ge0$ and $u$ a real smooth allowed
solution changes the Gaussian state while leaving the commutator and the
singularity unchanged. Appendix~\ref{app:C} gives the distributional test.

There is a stronger obstruction than the fact that one particular state is
entangled. In the vacuum representation of standard continuum free fields,
local von Neumann algebras for ordinary bounded open diamonds are type III
\cite{Haag}. In particular they are not full operator algebras
$\mathcal B(\mathcal H_1)$ of a Hilbert-space factor: the latter have
minimal projections and are type I. Thus a sharp spatial division does not
in general give a physical tensor decomposition in which each local algebra
is the full algebra of its factor, and a partial trace of a density matrix
is not the definition of restriction to a local algebra. Restriction is
simply evaluation of $\varpi$ on its subalgebra. An abstract Hilbert-space
isomorphism, or the tensor product for finitely many lattice sites, does
not resolve the operator-algebra obstruction. With a separating collar,
models satisfying the split property admit an intermediate type I factor;
shrinking the collar is not a proof of sharp factorization. Nor does a
spacelike division of initial data make successive \emph{time} intervals
independent quantum systems.

An algebra isomorphism always transports states by composition. To identify
two \emph{specified} cyclic state representations, one must also compare
the states: for centered Gaussian states transported by $S$, this means
$\mu_{\rm sew}(v,w)=\mu_M(Sv,Sw)$. If it holds, mapping the cyclic vectors
and the vectors generated from them preserves inner products and gives the
corresponding unitary after completion. A scalar commutator or a matched
classical flux alone does not give this unitary.

\begin{exercise}
Keep the two local covariances $\nu 1_2$ of the oscillator example, but
change the cross block to $\operatorname{diag}(c,-c)$. How large can
$|c|$ be at fixed $\nu\ge1/2$? Can a local pair with vacuum covariance
$\nu=1/2$ have any such correlation?
\end{exercise}
\noindent\emph{Hint.} In the sum pair the displacement variance is
$\nu+c$ while the momentum variance is $\nu-c$; in the difference pair
they are interchanged. Apply Gaussian positivity to each entire pair.
